
\subsubsection{Contamination Detection Methods}

\textbf{N-gram overlap detection.} The GPT-3 decontamination methodology \citep{brown2020language} established 13-gram matching as the standard for identifying verbatim overlap between training and evaluation data. Yang et al. [2023] extended this analysis to RedPajama, finding 8-18\% overlap with HumanEval and demonstrating that simple paraphrasing bypasses n-gram detection. These methods detect contamination presence but do not quantify performance impact.

\textbf{Membership inference attacks.} MIA-based approaches [Carlini et al., 2021; Shi et al., 2024] determine whether specific examples appeared in training data by analyzing model behavior signals: loss, perplexity, confidence distributions. The MIMIR benchmark \citep{duan2024mimir} standardizes these approaches for LLMs. However, MIA methods answer whether data was seen, not how much seeing it affected performance.

\textbf{Output distribution analysis.} Dong et al. [2024] proposed CDD (Contamination Detection via Distribution) to identify contaminated samples through output distribution peakedness, paired with TED (Test Data Deviation) for mitigation. Choi et al. [2025] introduced Kernel Divergence Score for dataset-level contamination detection via fine-tuning sensitivity. These methods remain focused on contamination presence rather than performance correlation.

\subsubsection{Benchmark Reliability}

\textbf{Contamination-resistant benchmarks.} Recent work has developed frequently-updated benchmarks resistant to training data leakage. LiveBench \citep{white2024livebench} provides monthly-updated questions with objective ground-truth scoring. LessLeak-Bench \citep{zhou2025lessleak} covers 83 software engineering benchmarks with automated anti-leakage measures. These efforts sidestep contamination rather than quantifying its effects.

\textbf{Benchmark analysis.} Sainz et al. [2023] documented contamination levels across major benchmarks and argued for community detection efforts. Their work establishes contamination as widespread but stops at detection.

\subsubsection{Learning Dynamics and Memorization}

\textbf{Training dynamics analysis.} Research on learning dynamics [Swayamdipta et al., 2020; Toneva et al., 2019] reveals that models learn different examples at different rates. This work inspired the checkpoint-gradient approach---if memorization accumulates during training, contamination effects should correlate with training progress.

\textbf{Memorization studies.} Carlini et al. [2023] demonstrated that LLMs memorize substantial training content, with extractable memorization scaling with model size and data repetition. Biderman et al. [2023] provided training dynamics for the Pythia model family with documented checkpoint availability, enabling checkpoint-gradient methodology.

\subsubsection{Position}

Existing work establishes that contamination exists, can be detected, and that models memorize training content. What remains missing is the quantitative bridge: given a certain contamination level, how much does the benchmark score inflate? This work addresses this gap through checkpoint-gradient analysis with capability detrending.

